% Section 1 of the CIRP CMS paper: introduction. Draft of 3 Nov 2026.
%
% Written from docs/paper-outline.md (section 1) and the paper's abstract on Overleaf (title
% "Asset Administration Shell-driven reconfiguration of skill-based IEC 61499 devices"), in its terms:
% product, process and resource AAS; translation rules; the three triggers; parameter, flow and
% structural changes. It claims only what is built (see the outline's notes on the abstract). Budget:
% about 750 words; with section 2 it fills the first two pages (page 1 also holds the title, the
% abstract and the keywords). Citation keys are those of docs/paper-references.bib.

\section{Introduction}
\label{sec:introduction}

% P1: the setting, a change that needs new behaviour, and what IEC 61499 and the AAS promise
Pharmaceutical production is moving from large batches on dedicated lines to small batches and
personalised products, made on lines assembled from modules~\cite{schwoerer2025modular}. Changeovers
between products become frequent, and a new product may need behaviour that no module has yet. A vial
is filled with the needle at a fixed depth; a prefilled syringe is filled from the bottom of its
barrel, the needle rising while the pump doses. IEC~61499 suits such changes: its management commands
create, connect and parameterise function blocks in a running device~\cite{iec61499}, without stopping
the parts that do not change~\cite{rooker2007zero}. Each module also comes with an Asset
Administration Shell (AAS), the standardised digital twin that describes what the module is, what it
can do and how it is called. Together they promise production that is both interoperable and
reconfigurable. How the content of the AAS can drive the reconfiguration of the control program,
however, is open.

% P2: where the state of the art leaves it
Skill-based engineering has made a module's functions explicit: a capability states what a resource
can do and a skill how it is done~\cite{koecher2023reference,css2022}, and the skills of a station are
composed from the primitive skills of its devices~\cite{schou2018skill}. The AAS describes both, in
standard submodels such as the Capability Description~\cite{idta02020}. Yet the AAS remains a
description of a program made elsewhere. Compositions of skills are generated at engineering time and
deployed as new code~\cite{dorofeev2021orchestrator}, or executed by an orchestrator outside the
controller~\cite{jensen2026service,sidorenko2024composition}. Planners derive process plans from AAS
models, but they sequence skills that the resources already
have~\cite{bernhard2024leverage,nabizada2026capability,schwoerer2026fond}. A change of behaviour thus
still means the vendor's engineering project, an IDE and a new deployment, while the integrator who
runs the line typically has only the running program and the AAS.

% P3: our earlier work, and the step this paper takes % cut first: down to its last sentence
In earlier work we orchestrated such modules from their AASs, with behaviour trees over a unified
namespace~\cite{jensen2026service}, and planned line-level policies from distributed AAS
descriptions~\cite{schwoerer2026fond}. Both took a module's skills as given. This paper turns to the
module itself: how a new skill is made from the AAS, and how it reaches the running program.

% P4: the approach
We make the AAS the desired state of the module's IEC~61499 program, in a framework that follows the
product--process--resource paradigm. The product's AAS states the processes the product requires and
their parameters. The process AAS binds the product's procedure, a block-structured flow drawn in
BPMN, to the line's modules and their skills. The module's resource AAS describes its skills, their
interfaces and which of them its program carries out. The program follows a fixed design in which a
composite skill is data rather than code. Its shell is the same in every module: an occupation that
lets one client at a time command the module, and a PackML state machine for the module. Primitive
skills are function block types, implemented once and never changed; a composite skill is a
subapplication of their instances; and every skill has the same state machine. Translation rules map
a primitive to its block type, a composite's flow to the event and data connections of its
subapplication, and its parameters to values written into it. The rules infer all wiring from the
flow alone, so the AAS needs to say nothing about blocks.

% P5: the two routes, the check and the one path for every change
A new composite enters the AAS in one of two ways. An engineer handcrafts it in a skill editor, in the
same flow elements as a product's procedure. Or a process step requires a capability that the module
does not offer, and the module plans its primitives into an order that reaches it, from the pre- and
postconditions of their capabilities. Every change then follows one path, whether it comes from a
product changeover, an edit of the module's skills or an operator's edit of a flow. It is checked
against the primitives' preconditions and effects and against the product's requirement, classified
as a parameter, flow or structural change, applied online through management commands where its class
allows it, verified by reading the program back and comparing it with the AAS, and recorded in the
resource AAS.

% P6: contributions
The contributions are:
\begin{itemize}
  \item a module architecture for IEC~61499 in which composite skills are networks of instances, with
    rules that map a skill in the AAS to the program and the program back to the skill (C1);
  \item an AAS model that joins a product's required processes, the module's composite skills and its
    components' primitive skills, through capabilities derived from the skills (C2);
  \item two routes to a new skill, handcrafted or planned, both checked against the primitives'
    conditions and the product's requirement before anything is deployed (C3);
  \item the reconciliation of the running program with the AAS: compare, apply online over the
    management protocol, verify by reading back, and record (C4).
\end{itemize}
The approach is demonstrated with Eclipse 4diac~\cite{strasser2008framework} on a pharmaceutical
filling module, with product variants that differ in fill volume, number of doses and inspection by
weighing, and with the changeover from vials to prefilled syringes.

% P7: structure % cut first
Section~\ref{sec:related} places the work among skill composition on controllers, the AAS as a source
for control, online reconfiguration of IEC~61499 and planning from capabilities.
Section~\ref{sec:method} presents the architecture and the method, Section~\ref{sec:evaluation} the
evaluation, and Section~\ref{sec:conclusion} concludes.
